
\section{4. Method}

\subsection{4.1 Overview}

Adaptive Span Attention (ASA-Span) replaces the standard softmax attention 
with a \textbf{span-masked} variant:

1. Compute standard scores $S = QK^\top / \sqrt{d_h}$.
2. Apply causal, padding, and position-bias masks.
3. Apply the \textbf{per-head soft ramp} $g_{z_h}(d)$.
4. Renormalize via softmax.
5. Multiply with $V$.

\textbf{Figure 4.1} illustrates the pipeline.

\begin{lstlisting}
    Q, K, V $\in$ R^{B $\times$ H $\times$ L $\times$ d_h}
             |
              
    +------------------------------+
    |  S = Q K^T / $\sqrt{}$d_h            |
    +--------------+---------------+
                   |
                    
    +------------------------------+
    |  Apply masks (causal, pad)   |
    +--------------+---------------+
                   |
                    
    +------------------------------+
    |  Per-head ramp:              |
    |  g_z(d) = clamp((R+z-d)/R)   |
    |  S  = S + log(g_z + $\epsilon$)        |
    +--------------+---------------+
                   |
                    
    +------------------------------+
    |  P = softmax(S )              |
    +--------------+---------------+
                   |
                    
              O = P   V
\end{lstlisting}

\subsection{4.2 Learnable Span}

\subsubsection{4.2.1 Parameterization}

For each of the $H$ heads, we introduce a scalar logit $\ell_h \in \mathbb{R}$. 
The span is:

\[
z_h = \sigma(\ell_h) \cdot S_\text{max}
\tag{4.1}
\]

where $\sigma$ is the sigmoid and $S_\text{max}$ is the maximum span 
(hyperparameter).

\subsubsection{4.2.2 Initialization}

We initialize $\ell_h$ to a value that makes $z_h \approx S_\text{max}$:

\[
\ell_h^{(0)} = \log \frac{0.982}{1 - 0.982} \approx 4.0
\tag{4.2}
\]

This gives $\sigma(4.0) \approx 0.982$, so $z_h \approx 0.982 \cdot S_\text{max}$. 
The span starts \textbf{nearly fully open} and shrinks during training via the 
L1 penalty.

\subsubsection{4.2.3 Why Start Open?}

Starting with a small span would immediately hurt accuracy. Starting open 
lets the model decide which spans to shrink, based on task loss.

\subsection{4.3 Soft Ramp Mask}

\subsubsection{4.3.1 Definition}

Given query position $i$ and key position $j$, define the distance:

\[
d_{ij} = \begin{cases} |i - j| & \text{if bidirectional} \\ \max(0, i - j) & \text{if causal} \end{cases}
\tag{4.3}
\]

For causal models, keys after the query ($j > i$) are masked separately; 
the ramp only applies to \textbf{past} keys.

The ramp mask is:

\[
g_{z_h}(d) = \text{clamp}\!\left(\frac{R + z_h - d}{R}, 0, 1\right)
\tag{4.4}
\]

\subsubsection{4.3.2 Properties}

\begin{itemize}
\item \textbf{Continuous:} $g_z$ is continuous in $z$ and $d$.
\end{itemize}
\begin{itemize}
\item \textbf{Piecewise-linear:} Three regimes (full, linear decay, zero).
\end{itemize}
\begin{itemize}
\item \textbf{Differentiable in $z$:}
\end{itemize}
\[
  \frac{\partial g_z(d)}{\partial z} = \begin{cases} 0 & d < z \\ 1/R & z \le d \le z + R \\ 0 & d > z + R \end{cases}
  \tag{4.5}
\]
\begin{itemize}
\item \textbf{Bounded:} $g_z \in [0, 1]$.
\end{itemize}

\subsubsection{4.3.3 Numerical Stability}

We compute the ramp in FP32 and apply in log space:

\[
\tilde{S}_{ij} = S_{ij} + \log(g_{z_h}(d_{ij}) + \epsilon)
\tag{4.6}
\]

with $\epsilon = 10^{-8}$. Then standard softmax.

\subsection{4.4 Renormalization}

After masking, the attention row does not sum to 1. We renormalize:

\[
P_{ij} = \frac{\exp(\tilde{S}_{ij})}{\sum_j \exp(\tilde{S}_{ij})}
\tag{4.7}
\]

This is equivalent to applying the ramp as a \textbf{multiplicative factor} before 
softmax:

\[
P_{ij} = \frac{g_{z_h}(d_{ij}) \cdot \exp(S_{ij})}{\sum_j g_{z_h}(d_{ij}) \cdot \exp(S_{ij})}
\tag{4.8}
\]

\subsection{4.5 Fast Path Optimization}

If the ramp has no effect (all heads have $z_h + R \geq L$), we skip ramp 
computation:

\begin{lstlisting}
if not ((span + self.ramp_size) >= N_k).all():
    # compute ramp
    ...
\end{lstlisting}

This avoids the $\mathcal{O}(H L^2)$ ramp allocation when unnecessary.

\subsection{4.6 L1 Span Penalty}

\subsubsection{4.6.1 Formulation}

We add to the loss:

\[
\mathcal{L}_\text{total} = \mathcal{L}_\text{task} + \lambda \sum_{h=1}^H z_h
\tag{4.9}
\]

\subsubsection{4.6.2 Effect}

The gradient w.r.t. $\ell_h$:

\[
\frac{\partial \mathcal{L}_\text{span}}{\partial \ell_h} = \lambda \cdot S_\text{max} \cdot \sigma(\ell_h)(1 - \sigma(\ell_h))
\tag{4.10}
\]

always negative (pushes $z_h$ toward 0). The task loss pushes $z_h$ up 
when the model needs long-range context. The equilibrium depends on 
$\lambda$.

\subsubsection{4.6.3 Recommended Values}

\begin{table}[H]
\centering
\begin{tabular}{ll}
\toprule
\textbf{$\lambda$} & \textbf{Effect} \\
\midrule
0.0 & No compression; spans stay open \\
0.001 & Mild compression \\
0.01 & Balanced (recommended) \\
0.1 & Aggressive; likely accuracy drop \\
\bottomrule
\end{tabular}
\end{table}

We use $\lambda = {{LAMBDA}}$ by default.

\subsection{4.7 Complete Forward Pass}

\textbf{Algorithm 1:} Adaptive Span Attention forward
\begin{lstlisting}
Input:  Q, K, V $\in$ R^{B $\times$ H $\times$ L $\times$ d_h}
        span_logits   $\in$ R^H
        max_span S_max, ramp_size R
        attention_mask, is_causal
Output: context $\in$ R^{B $\times$ H $\times$ L $\times$ d_h}

1:  span $\leftarrow$ sigmoid( ) * S_max                # (H,)
2:
3:  S $\leftarrow$ Q @ K^T * scale                      # (B, H, L, L)
4:
5:  S $\leftarrow$ apply_attention_mask(S, attention_mask)
6:  if is_causal:
7:      S $\leftarrow$ S + causal_mask
8:
9:  # Apply ramp mask
10: if (span + R).min() < L:
11:     d_q $\leftarrow$ arange(L)[:, None]
12:     d_k $\leftarrow$ arange(L)[None, :]
13:     if is_causal:
14:         dist $\leftarrow$ clamp(d_q - d_k, min=0)
15:     else:
16:         dist $\leftarrow$ abs(d_q - d_k)
17:     span_b $\leftarrow$ span.view(1, H, 1, 1)
18:     ramp $\leftarrow$ clamp((R + span_b - dist) / R, 0, 1)
19:     S $\leftarrow$ S + log(ramp + $\epsilon$)
20: end if
21:
22: P $\leftarrow$ softmax(S)
23: P $\leftarrow$ dropout(P)
24: O $\leftarrow$ P @ V
25: return O
\end{lstlisting}

\subsection{4.8 Complexity Analysis}

\begin{table}[H]
\centering
\begin{tabular}{lll}
\toprule
\textbf{Component} & \textbf{Dense} & \textbf{ASA-Span} \\
\midrule
Score matrix & $\mathcal{O}(B H L^2)$ & $\mathcal{O}(B H L^2)$ \\
Ramp mask & -- & $\mathcal{O}(B H L^2)$ \\
Softmax & $\mathcal{O}(B H L^2)$ & $\mathcal{O}(B H L^2)$ \\
\bottomrule
\end{tabular}
\end{table}

\textbf{Note.} ASA-Span does \textbf{not reduce complexity}. The full score matrix is 
still computed. Its benefit is \textbf{interpretability} and \textbf{pruning} -- the 
model learns which spans matter. To realize speedup, one would apply a hard 
mask after training and use sparse kernels.

\subsection{4.9 Parameter Analysis}

\begin{table}[H]
\centering
\begin{tabular}{lll}
\toprule
\textbf{Component} & \textbf{Baseline} & \textbf{ASA-Span} \\
\midrule
$W_Q$ & $d H d_h$ & Same \\
$W_K$ & $d H_{kv} d_h$ & Same \\
$W_V$ & $d H_{kv} d_h$ & Same \\
$W_O$ & $H d_h d$ & Same \\
**Span logits** & 0 & **$H$** \\
**LoRA** & 0 & $2 r_L (\cdots)$ \\
\bottomrule
\end{tabular}
\end{table}

ASA-Span adds \textbf{$H$ parameters} (one logit per head). It does not reduce 
parameter count.

\subsection{4.10 Theoretical Notes}

\subsubsection{4.10.1 Approximation Quality}

For a head with span $z$ and ramp $R$, the approximation error is bounded by:

\[
\|P - \tilde{P}\|_1 \leq \sum_{(i,j) : |i-j| > z} P_{ij} + \epsilon_R
\tag{4.11}
\]

where $\epsilon_R$ is the ramp-induced error:

\[
\epsilon_R \leq \frac{R}{\sum_j g_z(d_{ij}) \exp(S_{ij})} \cdot \sum_j \exp(S_{ij})
\tag{4.12}
\]

\subsubsection{4.10.2 Head Span Heterogeneity}

Empirically, spans converge to a bimodal distribution:

\[
z_h^{(T)} \approx \begin{cases} 16 & \text{if head is local} \\ S_\text{max} & \text{if head is global} \end{cases}
\tag{4.13}
\]

\subsubsection{4.10.3 Equilibrium of L1 Penalty}

The equilibrium span satisfies:

\[
\frac{\partial \mathcal{L}_\text{task}}{\partial z_h} + \lambda = 0
\tag{4.14}
\]

The task loss derivative depends on the model, but in practice heads with 
small gradients converge to small spans.